% note1.tex --- Track M: a standalone note for complex analysts, with no prerequisites
% from the Arithmetic Landscapes series.  Commissioned by fable-5 at r192, audited at
% r205 (Theorem 2 raised to *proved* by an independent reading along a disjoint route),
% and green-lit for a manuscript at r207.
%
% Statuses are carried at every statement, as everywhere in this project.  The one
% deliberate open item -- the constant at s = 1 -- is stated as unsettled with its full
% record, because an honest unknown in a note is a gift to whoever closes it.
\documentclass[11pt,a4paper]{article}
\usepackage[T1]{fontenc}
\usepackage{lmodern}
\usepackage{amsmath,amssymb,amsthm}
\usepackage{mathtools}
\usepackage{booktabs}
\usepackage[margin=28mm]{geometry}
\usepackage{microtype}
\usepackage[colorlinks,citecolor=blue,linkcolor=blue,urlcolor=blue]{hyperref}
\usepackage{xcolor}

\theoremstyle{plain}
\newtheorem{theorem}{Theorem}[section]
\newtheorem{proposition}[theorem]{Proposition}
\newtheorem{lemma}[theorem]{Lemma}
\newtheorem{corollary}[theorem]{Corollary}
\theoremstyle{definition}
\newtheorem{definition}[theorem]{Definition}
\newtheorem{problem}[theorem]{Problem}
\newtheorem{conjecture}[theorem]{Conjecture}
\theoremstyle{remark}
\newtheorem{remark}[theorem]{Remark}

\newcommand{\STATUS}[1]{\textcolor{blue}{\textsf{[#1]}}}
\newcommand{\Xref}[3]{#3}
\newcommand{\Xlab}[2]{\textup{\texttt{#2}}}

\title{Two speeds at the boundary:\\
zeros of sums of conjugate sections of power series}
\author{Kentaro Amauchi\thanks{Independent researcher.
\href{mailto:amkn.sub03@gmail.com}{amkn.sub03@gmail.com},
\url{https://mathlab0911.github.io/}, ORCID \url{https://orcid.org/0009-0000-0890-4395}}}
\date{\today}

\begin{document}
\maketitle

\begin{abstract}
Let $w_0,w_1,\dots$ be non-negative reals, $G(z)=\sum_{j\ge0}w_jz^{\,j}$ and
$G_k(z)=\sum_{j<k}w_jz^{\,j}$ its sections. We study the zeros of
\[
  F_k(q)\;=\;1+G_k(2q)+G_k(2-2q),
\]
a sum of two sections under substitutions that are complex conjugates of one another on
the line $\operatorname{Re}q=\tfrac12$, and in particular whether they approach the
symmetry point $q=\tfrac12$ as $k\to\infty$.

Jentzsch's theorem describes the zeros of a \emph{single} section, and adding functions
can cancel zeros; we found no result covering sums of sections. The observation this note
is built on is that on the symmetry line $F_k$ is \emph{real-valued}, so its zeros there
are sign changes --- and a sign change cannot be cancelled. Two results follow. First,
for constant weights we obtain a closed form
$F_k(\tfrac12+it)=A+w\rho^{\,k}\sin(k\theta)/t$ with $A=1+2w_0-2w$, from which the zero
nearest the symmetry point exists and satisfies $t_1=\pi/2k\,(1+O(1/k))$, and is given
exactly by $t=\tfrac12\tan(n\pi/k)$ when $A=0$. Second, when the weights decay the rate
changes: for $w_j\asymp(j+1)^{-s}$ we find $t_1\asymp\sqrt{s\log k/2k}$, slower by a
factor $\asymp\sqrt{k\log k}$. The dividing line is the decay of $w_j$ and not the
convergence of $\sum_jw_j$ --- a distinction that our own data did not test until a
profile lying on the line was computed. The constant at $s=1$ is reported as unsettled,
with the two candidate laws we have refuted or failed to resolve.
\end{abstract}

\section{The question}\label{sec:question}

Throughout, $w_0,w_1,w_2,\dots$ are non-negative real numbers, not all zero,
\[
  G(z)=\sum_{j\ge0}w_j z^{\,j},\qquad
  G_k(z)=\sum_{j=0}^{k-1}w_j z^{\,j},
\]
and $R\in(0,\infty]$ denotes the radius of convergence of $G$. The object of study is
\begin{equation}\label{eq:F}
  F_k(q)\;=\;1+G_k(2q)+G_k(2-2q),
\end{equation}
a polynomial of degree at most $k-1$, invariant under $q\mapsto1-q$. We ask where its
zeros go as $k\to\infty$, and specifically whether they approach the symmetry point
$q=\tfrac12$.

\paragraph{Why this is not immediate.}
The zeros of a \emph{single} section are classical. \textbf{Jentzsch's theorem} (1914)
states that if $0<R<\infty$ then every point of $|z|=R$ is a limit point of the zeros of
the sections $G_k$; Szeg\H{o}'s refinement gives angular equidistribution along a
subsequence, and the result has been extended to Dirichlet, Kapteyn and Neumann series,
to analytic curves and to ultrametric fields. Applied to $G$ with $R=1$, it says the
zeros of $G_k$ accumulate at $z=1$, which is $q=\tfrac12$ under $z=2q$.

But \eqref{eq:F} is a \emph{sum} of two sections under affine substitutions, plus a
constant, and adding functions can cancel zeros: nothing in the single-section theory
survives the sum automatically. We searched for a result covering sums of sections and
found none; \S\ref{sec:lit} says exactly where we looked, so that a reader who knows
otherwise can correct us cheaply.

\paragraph{Where the object comes from, and why it does not matter here.}
$F_k$ arose elsewhere, as a one-parameter deformation of a quantity attached to a finite
set of integers, in which the $w_j$ are a normalised gap statistic. \textbf{Nothing below
uses that}, and the note is written so that it \emph{cannot}: every statement from here on
is about $w_j$, $G$ and $F_k$ as defined above, and nothing else is assumed, cited or
required.

\subsection{Terminology}\label{sec:terminology-note}

\emph{This note carries no vocabulary of its own, and there is therefore no terminology
table.} Everything below is stated in the standard language of power series: coefficients,
radius of convergence, sections, and zeros. The absence of a table is the point rather
than an omission --- a reader should be able to check every statement here without
acquiring a single term.


\section{The reduction: on the symmetry line the object is real}\label{sec:reduction}

\subsection{Why these two substitutions}\label{sec:whytwo}

A reader may reasonably ask why \eqref{eq:F} pairs $G_k(2q)$ with $G_k(2-2q)$ rather than
with some other companion. The answer is that this is the only affine pairing for which
the two arguments are complex conjugates along a line, which is what makes
\S\ref{sec:reduction} available.

Concretely, let $\varphi_1(q)=\alpha q+\beta$ and $\varphi_2(q)=\gamma q+\delta$ with
real coefficients. Then $\varphi_2(q)=\overline{\varphi_1(q)}$ for all $q$ on a vertical
line $\operatorname{Re}q=c$ exactly when $\gamma=-\alpha$ and
$\delta=\beta+2\alpha c$ --- that is, when $\varphi_2$ is $\varphi_1$ reflected in that
line. Taking $\alpha=2$, $\beta=0$, $c=\tfrac12$ gives $\varphi_2(q)=2-2q$, which is
\eqref{eq:F}.

\begin{remark}[what the pairing costs and buys]\label{rem:pairing}
\STATUS{proved; elementary.}
The reflection is what creates the difficulty and also what resolves it. It creates it
because $G_k(\varphi_1)+G_k(\varphi_2)$ is a genuine sum, to which the single-section
theory does not apply. It resolves it because on the mirror line the sum is
$2\operatorname{Re}G_k(\varphi_1)$, a real function --- and the zeros of a real function
of one variable are visible by the intermediate value theorem, with no complex analysis
at all.
\end{remark}

\emph{This is the note in one paragraph.} The rest is the work of finding, for each class
of weights, a place where that real function is negative.



\begin{proposition}[reduction]\label{prop:real}
\STATUS{proved. Two independent derivations along disjoint routes (r203, r205);
verified numerically against direct summation at $80$ digits, worst agreement
$50$ digits, \texttt{trackm\_r203}.}
Let the $w_j$ be real. Then for every real $t$,
\[
  F_k\!\left(\tfrac12+it\right)\;=\;1+2\operatorname{Re}G_k(1+2it)
  \;=\;1+2\sum_{j<k}w_j\,\rho^{\,j}\cos(j\theta),
\]
where $\rho=|1+2it|=\sqrt{1+4t^{2}}$ and $\theta=\arg(1+2it)=\arctan 2t$. In particular
$F_k$ is real-valued on $\operatorname{Re}q=\tfrac12$, and every sign change of
$t\mapsto F_k(\tfrac12+it)$ is a zero of $F_k$ on that line.
\end{proposition}

\begin{proof}
With $q=\tfrac12+it$ we have $2q=1+2it$ and $2-2q=1-2it=\overline{2q}$. Since the $w_j$
are real, $G_k(\bar z)=\overline{G_k(z)}$, so $G_k(2q)+G_k(2-2q)=2\operatorname{Re}
G_k(1+2it)$. Writing $1+2it=\rho e^{i\theta}$ gives the sum. \qed
\end{proof}

Two remarks, and the second is the whole method of this note.

\begin{remark}[the shape of the sum]\label{rem:shape}
\STATUS{immediate from Proposition~\ref{prop:real}.}
The displayed sum is a cosine series with \emph{growing} amplitudes: $\rho>1$ for
$t\ne0$, so the coefficients $w_j\rho^{\,j}$ may increase even when the $w_j$ decrease.
Every estimate below is a competition between that growth and that decay.
\end{remark}

\begin{remark}[the cancellation problem disappears on the line]\label{rem:nocancel}
\STATUS{proved; it is the content of Proposition~\ref{prop:real}.}
The obstruction to transporting Jentzsch is that two sections may cancel each other's
zeros. On $\operatorname{Re}q=\tfrac12$ the object is a real function of one real
variable, and \textbf{a sign change of a real function cannot be cancelled by anything}.
So the on-line zeros are exactly the sign changes, and their existence is a question
about a real function --- which is why the results below are elementary where the
general problem is not.
\end{remark}

\begin{remark}[what the reduction does \emph{not} say]\label{rem:notnearest}
\STATUS{measured; the counterexample is explicit, \texttt{pinch\_r202}.}
Proposition~\ref{prop:real} says the line is where $F_k$ is real. \textbf{It does not say
the zero nearest $q=\tfrac12$ lies on it.} For $w_j=(c/2)^{j}$ with $c<2$ --- so $R>1$
--- it does not: at $c=1$, $k=70$ the nearest \emph{on-line} zero is at distance
$0.872$ while the nearest zero is at $0.504$, certified by the argument principle. The
nearest zero lies on the line in the boundary regime $R=1$, where it was confirmed by an
argument-principle count for $r\le0.25$, and this note asserts it only there.

\emph{We record the point because we imported the on-line assumption across the two
regimes ourselves, silently, and a winding-number certificate is what caught it.}
\end{remark}

\section{Constant weights: a closed form, and the pinch}\label{sec:const}

\begin{theorem}[constant weights]\label{thm:const}
\STATUS{proved. Written proof below; independently re-derived along a disjoint route by
a second reader who did not see it (r205); consequences \textup{(a)--(e)} and the three
constants of \S\ref{sec:instances} re-derived in that reading. Numerically verified,
\texttt{trackm\_r203}.}
Let $w_0\ge0$ and $w_j=w>0$ for $1\le j\le k-1$, and put
\[
  A\;:=\;1+2w_0-2w,\qquad \rho=\sqrt{1+4t^{2}},\qquad \theta=\arctan2t .
\]
Then for every $t>0$,
\begin{equation}\label{eq:closed}
  F_k\!\left(\tfrac12+it\right)\;=\;A+w\,\rho^{\,k}\,\frac{\sin(k\theta)}{t},
\end{equation}
exactly. Consequently:
\begin{enumerate}
\item[\textup{(a)}] $F_k(\tfrac12+it)\to1+2w_0+2(k-1)w$ as $t\to0^{+}$, the value at the
      symmetry point;
\item[\textup{(b)}] if $A\ge0$ then $F_k(\tfrac12+it)>0$ for $0<\theta<\pi/k$: no zero
      before the first rung;
\item[\textup{(c)}] at $\theta=3\pi/2k$, that is $t=\tfrac12\tan(3\pi/2k)$, one has
      $\sin(k\theta)=-1$ and hence
      $F_k\le A-w\,\tfrac{4k}{3\pi}\bigl(1+o(1)\bigr)<0$ for all large $k$;
\item[\textup{(d)}] therefore $F_k$ has a zero on the line with
      \[
        \tfrac12\tan(\pi/k)\;<\;t_1\;\le\;\tfrac12\tan(3\pi/2k),
        \qquad\text{so}\qquad t_1=\frac{\pi}{2k}\bigl(1+O(1/k)\bigr);
      \]
\item[\textup{(e)}] if $A=0$ the zero set on the line is \emph{exactly}
      $t=\tfrac12\tan(n\pi/k)$, $n=1,2,\dots,\lfloor(k-1)/2\rfloor$.
\end{enumerate}
\end{theorem}

\begin{proof}
The sum is geometric:
$G_k(z)=w_0+w\sum_{j=1}^{k-1}z^{\,j}=w_0+w\,(z^{k}-z)/(z-1)$, and on the line
$z-1=2it$, so
\[
  \operatorname{Re}\frac{z^{k}-z}{2it}
  =\frac{\operatorname{Im}(z^{k}-z)}{2t}
  =\frac{\rho^{k}\sin(k\theta)-2t}{2t}
  =\frac{\rho^{k}\sin(k\theta)}{2t}-1 .
\]
Substituting into Proposition~\ref{prop:real} gives
$F_k=1+2w_0+2w\bigl[\rho^{k}\sin(k\theta)/2t-1\bigr]$, which is \eqref{eq:closed}.

For (a), $\theta=2t+O(t^{3})$ and $\rho^{k}\to1$ as $t\to0$, so $\sin(k\theta)/t\to2k$
and $F_k\to A+2kw=1+2w_0+2(k-1)w$. For (b), $\sin(k\theta)>0$ on $0<\theta<\pi/k$ while
$w,\rho^{k},t>0$, so $F_k>A\ge0$. For (c), $t=\tfrac12\tan(3\pi/2k)\le(3\pi/4k)(1+O(k^{-2}))$
and $\rho^{k}\ge1$, so $w\rho^{k}/t\ge 4kw/3\pi\,(1+o(1))$, which exceeds the constant
$A$ once $k>3\pi A/4w$. Then (b) and (c) give opposite signs at the two ends of
$[\tfrac12\tan(\pi/k),\tfrac12\tan(3\pi/2k)]$ and $F_k$ is continuous there, so it
vanishes inside; both endpoints are $\tfrac{\pi}{2k}(1+O(1/k))$, which is (d). For (e),
$A=0$ leaves $w\rho^{k}\sin(k\theta)/t$, and $w,\rho^{k},t$ are non-zero for $t>0$. \qed
\end{proof}

\subsection{Three instances, and why one is special}\label{sec:instances}

The families that motivated this note are all instances of Theorem~\ref{thm:const},
distinguished only by the value of $A$.

\begin{center}
\begin{tabular}{lccc l}
\toprule
weights & $w_0$ & $w$ & $A$ & zero set on the line\\
\midrule
$w_0=0$, $w_j=\tfrac12$ & $0$ & $\tfrac12$ & $\mathbf{0}$
  & \emph{exactly} $t=\tfrac12\tan(n\pi/k)$\\
$w_0=1$, $w_j=\tfrac12$ & $1$ & $\tfrac12$ & $2$
  & bracketed, $t_1=\tfrac{\pi}{2k}(1+O(1/k))$\\
$w_0=1$, $w_j=1$ & $1$ & $1$ & $1$ & the same\\
\bottomrule
\end{tabular}
\end{center}

\begin{remark}[one identity, three constants]\label{rem:onethree}
\STATUS{proved, being Theorem~\ref{thm:const} evaluated three times.}
These three had been measured, cited and counted separately as evidence before the
general case was written down. They are one identity with one constant, and the row
admitting a closed form is exactly the row with $A=0$, i.e.\ the row whose $w_0$
vanishes.

\emph{Separately computed instances of one mechanism look like independent evidence and
are not: the count was a count of computations, not of cases.} Stating the general
member also \emph{explains} which row is special, which no amount of further measurement
would have done.
\end{remark}

\subsection{The $A=0$ family, written out}\label{sec:worked}

The row with $A=0$ is worth seeing in full, because it is the one case in this note where
the answer is a formula rather than a bracket. Take $w_0=0$ and $w_j=\tfrac12$ for
$1\le j\le k-1$. Then \eqref{eq:closed} reads
\[
  F_k\!\left(\tfrac12+it\right)\;=\;\frac{\rho^{\,k}\sin(k\theta)}{2t},
\]
so on the line $F_k$ vanishes precisely where $\sin(k\theta)=0$, i.e.\ at
$\theta=n\pi/k$, i.e.\ at $t=\tfrac12\tan(n\pi/k)$. The zeros therefore sit at
\[
  q\;=\;\tfrac12+\tfrac{i}{2}\tan\!\left(\frac{n\pi}{k}\right),
  \qquad n=1,\dots,\left\lfloor\tfrac{k-1}{2}\right\rfloor,
\]
and nowhere else on the line. The first three, at $k=32$:

\begin{center}
\begin{tabular}{crr}
\toprule
$n$ & $t_n=\tfrac12\tan(n\pi/32)$ & $k\,t_n/\pi$\\
\midrule
$1$ & $0.0492457016785821265$ & $0.5016126$\\
$2$ & $0.0994561836898290035$ & $1.013052$\\
$3$ & $0.151673341803671196$  & $1.544932$\\
\bottomrule
\end{tabular}
\end{center}

\STATUS{proved (Theorem~\ref{thm:const}(e)); the table is the formula evaluated, and
\textbf{every entry is confirmed to be a zero of the direct sum to sixty digits} as an
instrument control, \texttt{note1tab\_r208}. The first entry is reproduced independently
in three further experiments, \texttt{pinch\_r202}, \texttt{trackm\_r203},
\texttt{lambda\_r206}.}

\emph{Two features of this list are used later.} The spacing is uniform in $\theta$ ---
one rung per $\pi/k$ --- and $k\,t_n/\pi\to n/2$, so in $t$ the rungs are asymptotically
evenly spaced as well. Both facts fail as soon as the weights decay
(Remark~\ref{rem:s1notes}(i)), and the failure is not a small perturbation.

\begin{corollary}[the pinch at $R=1$, for eventually constant weights]\label{cor:pinch}
\STATUS{proved. \emph{Not} immediate from Theorem~\ref{thm:const}: that theorem admits
one exceptional weight $w_0$, whereas ``eventually constant'' admits finitely many. The
two lines that close the gap are given.}
Suppose $w_j=w>0$ for all $j\ge j_0$, with $w_0,\dots,w_{j_0-1}\ge0$ arbitrary. Then
$R=1$, and $F_k$ has a zero on $\operatorname{Re}q=\tfrac12$ with
$t_1=\dfrac{\pi}{2k}\bigl(1+O(1/k)\bigr)$.
\end{corollary}

\begin{proof}
Write $w_j=w+v_j$ for $j\ge1$, so $v_j=0$ once $j\ge j_0$, and put
$B:=2\sum_{1\le j<j_0}|v_j|$. By Proposition~\ref{prop:real},
\[
  F_k\!\left(\tfrac12+it\right)
  \;=\;\underbrace{A+w\rho^{\,k}\frac{\sin(k\theta)}{t}}_{\text{Theorem~\ref{thm:const}}}
  \;+\;2\sum_{1\le j<j_0}v_j\,\rho^{\,j}\cos(j\theta),
\]
and on $0\le\theta\le 3\pi/2k$ we have $1\le\rho^{\,j}\le\rho^{\,j_0}\to1$ for fixed
$j_0$, so the second sum is bounded by $B$ in absolute value for all large $k$,
\emph{uniformly in $k$}. It therefore changes nothing that grows.

At $\theta=3\pi/2k$ one has $\sin(k\theta)=-1$ and $t\le(3\pi/4k)(1+O(k^{-2}))$, so the
main term is at most $A-4kw/3\pi\,(1+o(1))$, which beats $A+B$ once
$k>3\pi(A+B)/4w$: the value is negative there, and (c) of
Theorem~\ref{thm:const} survives with $A$ replaced by $A+B$.

For the lower end, write $\theta=\pi/k-\varepsilon$ with $0<\varepsilon<\pi/k$. Then
$\sin(k\theta)=\sin(k\varepsilon)\ge \tfrac{2}{\pi}k\varepsilon$ and
$t=\tfrac12\tan\theta\le\pi/2k\,(1+O(k^{-2}))$, so the main term is at least
$(4/\pi^{2})wk^{2}\varepsilon\,(1+o(1))$, which exceeds $B$ as soon as
$\varepsilon\ge\pi^{2}B/(4wk^{2})$. Hence $F_k>0$ for
$\theta\le\pi/k-O(k^{-2})$, so $t_1\ge\tfrac12\tan\bigl(\pi/k-O(k^{-2})\bigr)
=\tfrac{\pi}{2k}\bigl(1-O(1/k)\bigr)$. Combining, $t_1=\tfrac{\pi}{2k}(1+O(1/k))$. \qed
\end{proof}

\begin{remark}[what the two lines buy]\label{rem:whybroader}
\STATUS{proved, being the corollary.}
The broad statement is the one Conjecture~\ref{conj:line}'s first branch needs to cite:
``eventually constant'' is a class, whereas Theorem~\ref{thm:const} as stated is a
family. The perturbation is bounded while the main term grows like $k$, which is why
finitely many arbitrary weights cost nothing --- \emph{and it is worth noticing that the
same estimate fails the moment infinitely many weights differ from $w$, which is exactly
where \S\ref{sec:decay} begins.}
\end{remark}

\section{Decaying weights: the second speed}\label{sec:decay}

Corollary~\ref{cor:pinch} settles the constant-weight case. When the weights decay the
pinch still occurs, and the rate is different.

\begin{proposition}[the scale, derived]\label{prop:scale}
\STATUS{derived, asymptotically and not to the standard of \S\ref{sec:const};
\textbf{measured} at four exponents and five sizes with no fitted constant,
\texttt{rate\_r200}, \texttt{divide\_r204}, \texttt{lambda\_r206}.}
Suppose $w_j\asymp(j+1)^{-s}$ with $s>0$. In Proposition~\ref{prop:real},
$\rho^{\,j}=e^{2jt^{2}(1+O(t^{2}))}$, so the last term of the sum carries
$w_k\rho^{\,k}\asymp k^{-s}e^{2kt^{2}}$. This is $o(1)$ below the scale
$2kt^{2}=s\log k$ and unbounded above it, so a zero cannot occur below that scale and
the first one is expected at it:
\[
  t_1\;\asymp\;\sqrt{\frac{s\log k}{2k}} .
\]
\emph{The scale is exact; the constant in front of it is not, and nothing here claims
one.}
\end{proposition}

The contrast with Corollary~\ref{cor:pinch} is a factor $\asymp\sqrt{k\log k}$:

\begin{center}
\begin{tabular}{rcccc}
\toprule
$k$ & $t_1$, $w_j=\tfrac12$ & $t_1$, $w_j=(j+1)^{-1}$ & $k\,t_1$ (left) & $k\,t_1$ (right)\\
\midrule
$64$   & $0.0245634248847336271$  & $0.1861996814519841$  & $1.5720592$ & $11.91677961$\\
$128$  & $0.0122743110544627221$  & $0.1409993892900241$  & $1.5711118$ & $18.04792183$\\
$256$  & $0.00613623118978313762$ & $0.106767212545108$   & $1.5708752$ & $27.33240641$\\
$512$  & $0.00306800007881170098$ & $0.07234505270561657$ & $1.570816$  & $37.04066699$\\
\bottomrule
\end{tabular}
\end{center}

\noindent
The left column is $\tfrac12\tan(\pi/k)$, which Theorem~\ref{thm:const}(e) gives exactly;
the right is measured. The two right-hand columns are the diagnostic: $k\,t_1$
\emph{converges} to $\pi/2=1.5707963$ on the left and \emph{grows} on the right. That
divergence is the whole content of the two-regime statement --- the second family does
not merely pinch more slowly, it pinches at a different power of $k$.
\STATUS{left column proved, \texttt{trackm\_r203}; right column measured,
\texttt{divide\_r204}.}

\subsection{Where the dividing line is}\label{sec:line}

\begin{proposition}[the line is decay, not summability]\label{prop:line}
\STATUS{the refutation is \textbf{measured} and decisive, \texttt{divide\_r204}; the
positive statement is \textbf{conjectured}, see Conjecture~\ref{conj:line}.}
The dividing line between the two rates is whether $w_j\to0$, and \emph{not} whether
$\sum_jw_j$ converges. The harmonic profile $w_j=(j+1)^{-1}$ has $\sum_jw_j=\infty$ and
takes the $\sqrt{\;}$ rate: $k\,t_1$ measures
$11.9,\,18.0,\,27.3,\,37.0,\,55.6$ at $k=64,\dots,1024$, growing rather than approaching
$\pi/2=1.5707963$. The same holds for $w_j=\bigl((j+2)\log(j+2)\bigr)^{-1}$, which
diverges more slowly still ($k\,t_1$ reaching $62.2$).
\end{proposition}

\begin{remark}[how the wrong line survived, and what the mechanism said]\label{rem:whyline}
\STATUS{methodological; no mathematical content beyond Proposition~\ref{prop:line}.}
Every divergent-side profile we had computed had \emph{constant} weights and every
convergent-side one decayed like a power, so the data were equally consistent with a
different statement --- \emph{constant versus decaying} --- that had not been written
down. The two separate only on profiles that decay without summing, and none had been
computed. \emph{A dividing line is a claim about the profiles that lie on it, and
evidence from far to either side does not test it.}

Moreover the estimate in Proposition~\ref{prop:scale} uses $w_k$ and never uses
$\sum_jw_j$: summability was never one of its hypotheses. \emph{When a conjecture and
its own mechanism disagree about what the hypothesis is, the mechanism is the one that
was derived.}
\end{remark}

\subsection{The constant at $s=1$, unsettled}\label{sec:s1}

Proposition~\ref{prop:scale} names the scale but not the constant in it. Writing
\[
  \lambda_{\mathrm{eff}}(k,s)\;:=\;\frac{2k\,t_1^{2}}{\log k},
\]
the derived prediction is $\lambda_{\mathrm{eff}}\to s$. We report the state of that
question rather than a conclusion.

\begin{center}
\begin{tabular}{lll}
\toprule
candidate & evidence & verdict\\
\midrule
$\lambda_{\mathrm{eff}}\to s-\tfrac12$ & $|\lambda_{\mathrm{eff}}-(s-\tfrac12)|$ stays
  in $[0.34,0.48]$ and barely shrinks & \textbf{not refuted --- see
  Remark~\ref{rem:oneshape}}\\
$\lambda_{\mathrm{eff}}\to s$ & $|\lambda_{\mathrm{eff}}-s|$ reaches
  $0.16,\,0.14,\,0.099,\,0.093,\,0.018$ at $s=1,\tfrac32,2,3,4$ & consistent,
  \textbf{non-monotone}\\
correction of size $1/\log k$ & $(\lambda_{\mathrm{eff}}-s)\log k$ does not settle
  & \textbf{not resolved}\\
correction of size $\log\log k/\log k$ & nor does
  $(\lambda_{\mathrm{eff}}-s)\log k/\log\log k$ & \textbf{not resolved}\\
\bottomrule
\end{tabular}
\end{center}

\STATUS{measured over $k\le2048$ at five exponents, \texttt{lambda\_r206}; the decision
rule and the two candidate shapes were registered before the run.}

The measurements behind that table, in full, are:

\begin{center}
\begin{tabular}{rrrr}
\toprule
$k$ & $\lambda_{\mathrm{eff}}$, $s=1$ & $\lambda_{\mathrm{eff}}$, $s=2$
    & $\lambda_{\mathrm{eff}}$, $s=4$\\
\midrule
$64$   & $1.0670656$  & $2.8134049$ & $6.3869596$\\
$128$  & $1.0489407$  & $2.0877026$ & $5.3173568$\\
$256$  & $1.0525199$  & $2.0014877$ & $4.3314313$\\
$512$  & $0.85911175$ & $1.978608$  & $4.3134221$\\
$1024$ & $0.87254853$ & $1.86631$   & $4.071501$\\
$2048$ & $0.83659266$ & $1.9008311$ & $3.9822736$\\
\bottomrule
\end{tabular}
\end{center}

\STATUS{measured, \texttt{lambda\_r206}; $40$ working digits; the first zero located by
scanning upward from $t=0$.}

Each column is near its target and none is monotone, and the departures at the largest
$k$ are of the same size as those at $k=128$. \emph{We report the sequences rather than
a fitted constant}: a fit over these would produce a number the data do not support, and
the honest content is that the scale is right and the constant is not established.
Remark~\ref{rem:oneshape} explains \emph{why} these columns sit near $s$ while the
mechanism's own asymptote is $s-\tfrac12$, and in doing so withdraws our refutation of
that asymptote.

\begin{remark}[on the observable]\label{rem:observable}
\STATUS{methodological.}
The natural thing to print is $t_1$ divided by $\sqrt{s\log k/2k}$, and we did so for
four rounds. That ratio is $\sqrt{\lambda_{\mathrm{eff}}/s}$, and \textbf{a square root
halves relative error}: a $14\%$ departure in the quantity the derivation names appears
as $7\%$ on the page. \emph{Report the quantity the mechanism names; a monotone
transform is not a neutral change of units but a lens with a computable magnification.}
\end{remark}

\begin{remark}[two readings we can rule out]\label{rem:s1notes}
\STATUS{measured, \texttt{envelope\_r206b}.}
Two natural explanations for the non-monotonicity fail.
\begin{enumerate}
\item[(i)] \emph{Quantisation.} One might expect $t_1$ to be pinned to a ladder of
      troughs spaced $\pi/k$, as in Theorem~\ref{thm:const}(e), giving jitter
      $\pi/(kt_1)$. But when the weights decay the zeros are \emph{not} evenly spaced:
      the measured gaps run $0.17$ to $0.73$ in units of $\pi/k$, at $s=1$ and $s=2$
      alike. \textbf{The even ladder is a constant-weight phenomenon.}
\item[(ii)] \emph{An envelope threshold.} One might locate the first zero as the first
      $t$ with $2|G_k(1+2it)|=1$. No such $t$ exists: $2|G_k|-1$ equals $+11.2$ at
      $t=0$ for $s=1$, $k=1024$, since $|G_k(1)|$ is half the value at the symmetry
      point. The zero is produced by the \emph{phase} turning, not by the envelope
      growing.
\end{enumerate}
\end{remark}

\section{Two theorems, and the shape of the answer}\label{sec:twothm}

Since the first version of this note, the decaying case has stopped being a table of
measurements. Two statements are now proved, and together they fix the shape of the answer.

\begin{theorem}[every non-increasing profile]\label{thm:monotone}
\STATUS{proved. Two independent derivations along disjoint routes; the identity it rests on
was verified against direct summation to sixty digits, including at the zeros.}
Let $w_0\ge w_1\ge\dots\ge w_{k-1}\ge0$. If $\arctan(2t)\le\pi/k$ then $F_k(\tfrac12+it)\ge1$.
Consequently the first zero satisfies $t_1>\tfrac12\tan(\pi/k)$.
\end{theorem}

The proof is two lines from an exact identity obtained by summation by parts: writing
$D_j=w_{j-1}-w_j$ and $z=1+2it=\rho e^{i\theta}$,
\begin{equation}\label{eq:abel}
  F_k(\tfrac12+it)\;=\;1+\frac1t\Big[\,w_{k-1}\rho^k\sin(k\theta)
    +\sum_{j=1}^{k-1}D_j\,\rho^j\sin(j\theta)\,\Big].
\end{equation}
If the profile is non-increasing every $D_j$ is non-negative, and on $\theta\le\pi/k$ every
$j\theta$ with $j\le k$ lies in $(0,\pi]$, so every sine is non-negative and the bracket cannot
be negative. Theorem~\ref{thm:const} is the special case in which the sine series is empty.
\emph{Monotonicity is doing work and is not necessary}: an increasing ramp drives $F_k$ well
below $1$, while a bump and an alternating profile stay above it without being monotone.

\begin{theorem}[the second speed, for power weights]\label{thm:decayrate}
\STATUS{proved. Written proof and an independent re-derivation along a different route; the
two moment identities it uses are additionally replayed through the Lean kernel.}
Fix $s>1$ and let $w_j=(j+1)^{-s}$. Then
\[
  \frac{2k\,t_1(k)^2}{\log k}\;\longrightarrow\;s-\tfrac12\qquad(k\to\infty).
\]
\end{theorem}

The mechanism is visible in \eqref{eq:abel}. The coefficients are the weight
\emph{decrements}, which decay one power faster than the weights themselves, so the sine
series concentrates near $j\asymp1/2t$ rather than near $j\asymp k$; its limit is $2\zeta(s)$,
and the top term must supply the rest. The tail is controlled by cancellation and not by
absolute values --- the sum of $|D_j\rho^j|$ over the tail diverges for $1<s<2$ at this scale
--- which is possible because $D_j\rho^j$ has exactly one interior minimum, $D_j$ being
log-convex as the integral of a log-convex function over a unit interval.

\begin{remark}[what the two speeds are, and where they meet]\label{rem:shape}
\STATUS{$s>1$ \textbf{proved} (Theorem~\ref{thm:decayrate}); $s<1$ \textbf{measured},
decisively, the two candidate limits differing by more than twenty times the observable's own
resolution; $s=1$ \textbf{derived} by two independent routes.}
For $w_j=(j+1)^{-s}$ the effective exponent $\lambda=2kt_1^2/\log k$ appears to satisfy
\[
  \lambda_\infty(s)\;=\;\max\Big(\frac{s}{2},\;s-\frac12\Big),
\]
with the kink at $s=1$ --- exactly where $\sum_j w_j$ stops converging. Above the kink the
convergent weight sum sets the level and the top term alone places the zero; below it the sum
diverges and drags the zero outward. The constant governing $s<1$ is
$2^{s-1}\pi/(\Gamma(s)\cos(\pi s/2))$, whose surviving poles sit at the odd integers.
\end{remark}

\section{The general question}\label{sec:general}

\begin{problem}[does $R=1$ force the pinch?]\label{prob:general}
\STATUS{\textbf{proved} for constant weights (Corollary~\ref{cor:pinch});
\textbf{measured} for $w_j\asymp(j+1)^{-s}$ at four exponents; \textbf{open} in general.}
Let $w_j\ge0$ with $R=1$. Do the zeros of $F_k$ approach $q=\tfrac12$ as $k\to\infty$?
\end{problem}

\begin{conjecture}[two speeds]\label{conj:line}
\STATUS{conjectured; the constant-weight half is Corollary~\ref{cor:pinch} and is
proved, the decaying half is derived and measured, and the constant at $s=1$ is
unsettled (\S\ref{sec:s1}).}
Let $w_j\ge0$ with $R=1$. Then $F_k$ has a zero on $\operatorname{Re}q=\tfrac12$
approaching $q=\tfrac12$, at rate
\[
  t_1\sim\frac{\pi}{2k}\quad\text{if }w_j\not\to0,
  \qquad\qquad
  t_1\sim\sqrt{\frac{s\log k}{2k}}\quad\text{if }w_j\asymp(j+1)^{-s},\ s>0 .
\]
\emph{The two displayed classes are illustrative and do not exhaust $R=1$.} They omit,
for instance, profiles decaying more slowly than every power, such as
$w_j=1/\log(j+2)$. We expect such profiles on the second branch with a local exponent in
the role of $s$; that expectation is measured only at
$w_j=\bigl((j+2)\log(j+2)\bigr)^{-1}$ and is deliberately \textbf{not} part of the
conjecture.
\end{conjecture}

\paragraph{The missing analytic ingredient, named rather than left to be located.}
What is absent for the general case is the non-vanishing of the limit of $G$ on its
circle of convergence. Hurwitz's theorem gives non-accumulation on compact subsets of
the region of analyticity; accumulation \emph{on} the boundary needs that non-vanishing,
and it is a statement about $G$ alone.

\paragraph{A fingerprint of the sum, with its scope.}
Sections of a single geometric series have zeros at angular spacing $2\pi/k$; in the
constant-weight case ours sit at $\theta_n=n\pi/k$, twice as densely. That doubling is
what $\sin(k\theta)$ does and $e^{ik\theta}$ does not: taking a real part folds
$e^{ik\theta}$ into $\sin(k\theta)$, whose zeros are twice as dense. \textbf{So the
doubled ladder is exactly the part of the phenomenon the single-section theory does not
describe} --- and, by Remark~\ref{rem:s1notes}(i), it is a statement about profiles with
$w_j\not\to0$ and not a general one.
\STATUS{proved for constant weights (Theorem~\ref{thm:const}(e)); the scope restriction
is measured, \texttt{envelope\_r206b}.}

\section{What a proof of the decaying case would have to supply}\label{sec:whatproof}

Theorem~\ref{thm:const} works because the sum is geometric and collapses. When the
weights decay there is no such collapse, and we record what we tried, so that the next
attempt does not repeat it.

\paragraph{The naive balance is wrong.}
Near the first zero the sum in Proposition~\ref{prop:real} is dominated by $j$ close to
$k$, where the phases $j\theta$ turn rapidly. Integrating by parts once suggests the
oscillatory part has size $\asymp w_k\rho^{\,k}/\theta$ rather than $w_k\rho^{\,k}$, and
balancing \emph{that} against the constant term gives
$k^{-s}e^{2kt^{2}}\asymp t$, hence $\lambda_{\mathrm{eff}}\to s-\tfrac12$.

\begin{proposition}[the naive balance, and why it was wrongly refuted]\label{prop:notshalf}
\STATUS{\textbf{withdrawn as a refutation.} It was recorded as one on the grounds that
$|\lambda_{\mathrm{eff}}-(s-\tfrac12)|$ stays in $[0.34,0.48]$ over $k\le2048$ without
shrinking. Remark~\ref{rem:oneshape} shows that band is exactly what the balance
\emph{predicts} at those $k$. \texttt{lambda\_r206}, \texttt{rung0\_r211}.}
The balance $k^{-s}e^{2kt^{2}}\asymp t$ gives $\lambda_{\mathrm{eff}}\to s-\tfrac12$ with
an approach of size $(\log\log k)/(2\log k)$, which at $k=1024$ is $0.13966$. Over the
computable range it therefore places $\lambda_{\mathrm{eff}}$ near $s$, not near
$s-\tfrac12$, and the measurement cannot separate the two.
\end{proposition}

\begin{remark}[the two candidate shapes are one model at different $k$]\label{rem:oneshape}
\STATUS{the model is \textbf{derived}; its agreement with the measurement is
\textbf{measured} at $25$ pairs $(k,s)$ with no fitted parameter, worst error $9.9\%$,
\texttt{rung0\_r211}. The asymptote is a consequence of the model and is
\textbf{conjectured}.}
Split $G_k$ into a head and a tail. For small $j$ the terms contribute
$H:=\sum_{j<k}w_j$, real and positive; near $j=k$ the weights are nearly constant over
the range on which $|z|^{\,j}$ changes, so the tail is geometric,
\[
  \sum_{j\text{ near }k}w_jz^{\,j}\;\approx\;w_{k-1}\frac{z^{k}}{z-1},
  \qquad\text{modulus } T:=\frac{w_{k-1}\rho^{\,k}}{2t},
  \qquad\text{argument } k\theta-\frac{\pi}{2},
\]
the quarter turn coming from $\arg(z-1)^{-1}=-\pi/2$ exactly. Hence
$\operatorname{Re}G_k\approx H+T\sin(k\theta)$, and a zero of $F_k$ first becomes
possible when $T=H+\tfrac12$. \textbf{Solving that for $t$, with nothing fitted,
reproduces the measured first zero to within $9.9\%$ at every $(k,s)$ tested.}

Asymptotically the same equation reads $k^{\lambda-s}=2t(H+\tfrac12)$, so
\[
  \lambda_{\mathrm{eff}}-s\;=\;-\tfrac12+\frac{\log\log k}{2\log k}+O\!\left(\frac{1}{\log k}\right).
\]
\textbf{The two shapes registered in \S\ref{sec:s1} are therefore the same statement at
different $k$}: the correction term is $0.13966$ at $k=1024$ and decays like $1/\log k$,
so over any computable range $\lambda_{\mathrm{eff}}$ sits near $s$ while tending to
$s-\tfrac12$. \emph{Our own refutation of the second shape tested a limit that this model
says is approached far too slowly to be seen, and is withdrawn.}
\end{remark}

\paragraph{Two instruments that do not exist or do not apply.}
Remark~\ref{rem:s1notes} records both: the zeros are not a uniform ladder once the
weights decay, so no quantisation argument is available; and the first zero is not an
envelope crossing, because $2|G_k(1+2it)|$ already exceeds $1$ at $t=0$ by a wide margin.
\textbf{The zero is a phase event}: at the first zero, $\operatorname{Re}G_k=-\tfrac12$
while $|G_k|$ remains large, so $\cos(\arg G_k)=-1/(2|G_k|)$, which is close to $0$ from
below. A treatment in terms of $\arg G_k(1+2it)$ therefore has the right shape --- but a
linear approximation to that phase predicts a crossing at $t\asymp k^{-(2-s)}$, which is
not the observed scale, so any such treatment owes an account of how the amplitude
redistributes with $t$.
\STATUS{derived, at the level of a sketch; the linear-phase estimate is recorded because
it is wrong in a specific and informative way.}

\section{What is verified, and how}\label{sec:verified}

This note carries a status at every statement, and the statuses are meant literally.
Their evidence is as follows.

\begin{center}
\begin{tabular}{ll}
\toprule
status & what stands behind it here\\
\midrule
\emph{proved} & a written proof, plus an independent re-derivation along a disjoint
  route by a second reader who had not seen it\\
\emph{derived} & an asymptotic argument, stated as such; not to the standard of the
  proofs in \S\ref{sec:const}\\
\emph{measured} & a computation with a committed log, with the range stated; a range is
  not a limit\\
\emph{refuted} & a pre-registered falsifier fired\\
\emph{conjectured}, \emph{open} & neither of the above\\
\bottomrule
\end{tabular}
\end{center}

Two consequences of taking this seriously are visible in the text above.
Theorem~\ref{thm:const} says \emph{proved} because it was written by one reader and
re-derived from $G_k(z)=w_0+w(z-z^{k})/(1-z)$ by another who did not read the first
derivation; Proposition~\ref{prop:scale} says \emph{derived} and not \emph{proved},
although we believe it, because an asymptotic balance is not a proof.

\begin{remark}[the record of this note's own errors]\label{rem:ourerrors}
\STATUS{methodological; no mathematical content.}
Three statements that appear above in corrected form were previously published, under a
DOI, in incorrect form: an oscillator given as $\cos$ where the geometric sum gives
$\sin$ (the phase of $1/(z-1)$ is exactly $-\pi/2$, and a quarter turn moves a ladder by
half a rung); an equivalence ``$R\ge1$ if and only if $\sum_jw_j<\infty$'', which is false
exactly at $R=1$ where both cases occur; and the dividing line of
Conjecture~\ref{conj:line} given as summability rather than decay
(Remark~\ref{rem:whyline}). Each was corrected at its statement in the version where it
appeared, and named in that version's release notes; the record is public in the
repository identified in the disclosure below.

\emph{We record this because a note whose apparatus is part of its claim should say what
that apparatus has caught.} Two of the three were found not by a check but by building
something new and noticing that it contradicted a sentence written elsewhere.
\end{remark}

\section{Three questions}\label{sec:questions}

\begin{enumerate}
\item \textbf{Does a sum of two sections inherit boundary accumulation?} Precisely: for
      $w_j\ge0$ with $R=1$, does $F_k$ have zeros approaching $q=\tfrac12$?
      (Problem~\ref{prob:general}. Answered for $w_j\not\to0$;
      measured for $w_j\asymp(j+1)^{-s}$; open in general.)
\item \textbf{Is the dividing line exactly $w_j\to0$?} We have refuted summability as the
      criterion and verified the decay criterion on both sides, including on profiles
      that decay without summing --- but ``$w_j\to0$'' has not been shown to be
      \emph{sharp}. A profile with $w_j\to0$ taking the $\pi/2k$ rate would refute
      Conjecture~\ref{conj:line}.
\item \textbf{What is the constant at $s=1$?} The scale $\sqrt{s\log k/2k}$ is derived
      and measured at four exponents; at $s=1$ the constant is not established
      (\S\ref{sec:s1}). \emph{This is the cheapest of the three to attack and the one we
      most expect to be closed by someone else.}
\end{enumerate}

\section{Where we looked in the literature}\label{sec:lit}

So that the claim ``we found nothing covering sums of sections'' can be checked or
refuted cheaply, here is its extent. We consulted the classical statements of Jentzsch
(1914) and Szeg\H{o}, and surveyed the extensions of Jentzsch--Szeg\H{o} to Dirichlet
series, Kapteyn series, Neumann series, sections along analytic curves, and the
ultrametric setting. All of these concern the zeros of the sections of \emph{one}
series, or of one series composed with a map. We did not find a treatment of
\[
  \alpha\,G_k(\varphi_1(q))+\beta\,G_k(\varphi_2(q))+\gamma
\]
for affine $\varphi_1,\varphi_2$ --- in particular, of the conjugate pair that
\eqref{eq:F} produces on a line. \STATUS{a negative result about our search, not about
the literature; stated so that it can be overturned by one reference.}

\section{Numerical conventions}\label{sec:num}

Every number in this note is copied from a committed log, and every script named in a
status is in the repository with its log beside it. Computations use \texttt{mpmath} at
$30$--$120$ decimal digits, with the working precision printed; where a quantity is
compared across precisions, the number of agreeing digits is printed rather than
asserted. Zeros on the line are located by scanning $t$ upward \emph{from zero} and
bisecting the first sign change, so that ``first'' is a claim about $(0,t_{\max})$ and
not about a window; where a claim concerns zeros off the line as well, it is certified
by an argument-principle count on a circle about $q=\tfrac12$.

Each experiment registers its falsifiers before it runs, and at least one of them tests
the \emph{instrument} against a value proved elsewhere in this note --- typically
Theorem~\ref{thm:const}(e), whose zero set is exact. In the run that produced
\S\ref{sec:decay} the three falsifiers testing the law all passed while both instrument
controls failed, on an off-by-one in the weight index; the law-testing falsifiers alone
would have confirmed a wrong apparatus.

\section*{Tool and computational resource disclosure}
\addcontentsline{toc}{section}{Tool and computational resource disclosure}

This note was produced by one person using AI models as tools under his direction. Draft
text, proofs, code and numerical experiments were generated in collaboration with large
language models (Anthropic's Claude, models \texttt{opus-5} and \texttt{fable-5}); the
author directed the work, chose what to keep, and is responsible for the content.

Because work produced this way cannot be trusted on the author's word, the verification
apparatus is part of the artefact and is public: every experiment writes a committed
log, every number in this note is copied from one, statuses are carried at each
statement, and a failure ledger recording the project's own errors --- including three
corrections to statements that had already been published under a DOI --- is in the
repository. Theorem~\ref{thm:const} was written by one model and independently
re-derived along a disjoint route by the other before being labelled \emph{proved}.

Everything named in a status above --- every script, its log, the ledger, and the
twenty-one mechanical checks that gate each commit --- is at
\url{https://github.com/mathlab0911/arithmetic-landscapes}. \emph{A reader who wants to
disbelieve a number in this note should be able to find the line that produced it without
asking the author.}



\end{document}
